% Options for packages loaded elsewhere
\PassOptionsToPackage{unicode}{hyperref}
\PassOptionsToPackage{hyphens}{url}
\documentclass[
]{article}
\usepackage{xcolor}
\usepackage[margin=1in]{geometry}
\usepackage{amsmath,amssymb}
\setcounter{secnumdepth}{-\maxdimen} % remove section numbering
\usepackage{iftex}
\ifPDFTeX
  \usepackage[T1]{fontenc}
  \usepackage[utf8]{inputenc}
  \usepackage{textcomp} % provide euro and other symbols
\else % if luatex or xetex
  \usepackage{unicode-math} % this also loads fontspec
  \defaultfontfeatures{Scale=MatchLowercase}
  \defaultfontfeatures[\rmfamily]{Ligatures=TeX,Scale=1}
\fi
\usepackage{lmodern}
\ifPDFTeX\else
  % xetex/luatex font selection
\fi
% Use upquote if available, for straight quotes in verbatim environments
\IfFileExists{upquote.sty}{\usepackage{upquote}}{}
\IfFileExists{microtype.sty}{% use microtype if available
  \usepackage[]{microtype}
  \UseMicrotypeSet[protrusion]{basicmath} % disable protrusion for tt fonts
}{}
\makeatletter
\@ifundefined{KOMAClassName}{% if non-KOMA class
  \IfFileExists{parskip.sty}{%
    \usepackage{parskip}
  }{% else
    \setlength{\parindent}{0pt}
    \setlength{\parskip}{6pt plus 2pt minus 1pt}}
}{% if KOMA class
  \KOMAoptions{parskip=half}}
\makeatother
\usepackage{graphicx}
\makeatletter
\newsavebox\pandoc@box
\newcommand*\pandocbounded[1]{% scales image to fit in text height/width
  \sbox\pandoc@box{#1}%
  \Gscale@div\@tempa{\textheight}{\dimexpr\ht\pandoc@box+\dp\pandoc@box\relax}%
  \Gscale@div\@tempb{\linewidth}{\wd\pandoc@box}%
  \ifdim\@tempb\p@<\@tempa\p@\let\@tempa\@tempb\fi% select the smaller of both
  \ifdim\@tempa\p@<\p@\scalebox{\@tempa}{\usebox\pandoc@box}%
  \else\usebox{\pandoc@box}%
  \fi%
}
% Set default figure placement to htbp
\def\fps@figure{htbp}
\makeatother
\setlength{\emergencystretch}{3em} % prevent overfull lines
\providecommand{\tightlist}{%
  \setlength{\itemsep}{0pt}\setlength{\parskip}{0pt}}
\usepackage{bookmark}
\IfFileExists{xurl.sty}{\usepackage{xurl}}{} % add URL line breaks if available
\urlstyle{same}
\hypersetup{
  hidelinks,
  pdfcreator={LaTeX via pandoc}}

\author{}
\date{\vspace{-2.5em}}

\begin{document}

\section{Speaker Notes --- OQG Presentation (Zoe's
Section)}\label{speaker-notes-oqg-presentation-zoes-section}

Reference deck: \texttt{presentation/presentation\_d4.pptx}. My section
runs from slide 9 to the end. Total budget about 4 to 5 minutes. Format
follows my thesis speaker notes: short bullets, indented sub-bullets,
spoken aloud register.

\begin{center}\rule{0.5\linewidth}{0.5pt}\end{center}

\subsection{\texorpdfstring{Slide 9: Testing the Signal Against the
Underlying
\emph{(transition)}}{Slide 9: Testing the Signal Against the Underlying (transition)}}\label{slide-9-testing-the-signal-against-the-underlying-transition}

\textbf{Zoe takes over from Xavi, \textasciitilde30 sec}

\begin{itemize}
\tightlist
\item
  Now I take over from Xavi
\item
  A quick framing for how my piece fits in:

  \begin{itemize}
  \tightlist
  \item
    Xavi's work has been signal development and research

    \begin{itemize}
    \tightlist
    \item
      he built the Kalshi panel, showed it's calibrated, characterized
      the within-Kalshi signal
    \end{itemize}
  \item
    my piece is the next step

    \begin{itemize}
    \tightlist
    \item
      testing that signal against an external benchmark for prediction
    \item
      the underlying asset that every Kalshi contract resolves on
    \end{itemize}
  \end{itemize}
\item
  Every Kalshi contract resolves on BTC spot

  \begin{itemize}
  \tightlist
  \item
    so for Kalshi to be useful for trading the underlying
  \item
    it has to clear the bar of ``more than the spot tape already tells
    you''
  \end{itemize}
\item
  For the spot benchmark I pulled Coinbase BTC-USD 1-min candles

  \begin{itemize}
  \tightlist
  \item
    same 6-week window as Kalshi, joined by exact UTC minute
  \item
    31,658 matched minute-level observations
  \end{itemize}
\item
  The diagram on the right shows the three diagnostics I'm about to walk
  through

  \begin{itemize}
  \tightlist
  \item
    lead-lag cross-correlogram
  \item
    variance ratio test
  \item
    incremental information test
  \item
    each one motivated by the one before it
  \end{itemize}
\end{itemize}

\begin{center}\rule{0.5\linewidth}{0.5pt}\end{center}

\subsection{Slide 10: Spot Leads P(UP) by Two
Minutes}\label{slide-10-spot-leads-pup-by-two-minutes}

\textbf{\textasciitilde60 sec}

\begin{itemize}
\tightlist
\item
  First diagnostic is a timing question

  \begin{itemize}
  \tightlist
  \item
    at minute granularity, does P(UP) move before spot or after?
  \item
    if Kalshi is a leading sentiment signal, it should update before the
    underlying
  \item
    if it's downstream of spot, it should update after
  \end{itemize}
\item
  Method on this one is intuitive: Pearson cross-correlation at
  different time lags

  \begin{itemize}
  \tightlist
  \item
    correlate within-contract change in P(UP) at minute t
  \item
    with spot 1-min log return at minute t plus k
  \item
    sweep k from -10 to +10 minutes
  \item
    cluster bootstrap CIs over contracts (Cameron, Gelbach and Miller
    2008) to handle within-contract autocorrelation
  \end{itemize}
\item
  The result is the curve on screen
\item
  Peak is at lag k = -2 with

  \begin{itemize}
  \tightlist
  \item
    r = +0.510 for BTC
  \item
    r = +0.474 for ETH
  \end{itemize}
\item
  The contemporaneous correlation at lag zero is statistically zero

  \begin{itemize}
  \tightlist
  \item
    so this is not a same-minute relationship I am splitting
  \end{itemize}
\item
  Every lag at which Kalshi could lead spot, k ≥ 0

  \begin{itemize}
  \tightlist
  \item
    is statistically zero or negative
  \end{itemize}
\item
  The interpretation:

  \begin{itemize}
  \tightlist
  \item
    spot moves first
  \item
    P(UP) follows about two minutes later
  \end{itemize}
\item
  Some of the 2-min lag is mechanical

  \begin{itemize}
  \tightlist
  \item
    Coinbase candle aggregation contributes about 0.5 min
  \item
    Kalshi scraper averages mid-prices over the prior minute, another
    \textasciitilde0.5 min
  \item
    but at least 1 full minute of real reaction lag remains
  \end{itemize}
\item
  This 2-min lag was striking enough that I wanted a second test from a
  different angle

  \begin{itemize}
  \tightlist
  \item
    the cross-correlogram depends on the join between two separate
    panels
  \item
    if the same finding shows up internal to Kalshi alone, that is much
    stronger evidence
  \item
    which is the next slide
  \end{itemize}
\end{itemize}

\begin{center}\rule{0.5\linewidth}{0.5pt}\end{center}

\subsection{Slide 11: P(UP) Deviates from the Martingale
Null}\label{slide-11-pup-deviates-from-the-martingale-null}

\textbf{\textasciitilde70 sec (method + result combined)}

\begin{itemize}
\tightlist
\item
  Second diagnostic: forget the cross-series join for a minute

  \begin{itemize}
  \tightlist
  \item
    does P(UP) behave like a random walk on its own?
  \end{itemize}
\item
  The Lo and MacKinlay 1988 variance ratio test is the standard tool

  \begin{itemize}
  \tightlist
  \item
    source on the slide footer with DOI link
  \end{itemize}
\item
  Under the random-walk null, q-step changes have variance exactly q
  times the variance of 1-step changes

  \begin{itemize}
  \tightlist
  \item
    so the variance ratio VR(q) equals 1 at every horizon
  \item
    VR \textgreater{} 1: positive serial correlation, under-reaction,
    momentum
  \item
    VR \textless{} 1: mean reversion
  \end{itemize}
\item
  The Hurst exponent reads directly off the slope of log VR(q) vs log q

  \begin{itemize}
  \tightlist
  \item
    H = 0.5 is the random-walk reference
  \end{itemize}
\item
  I run this at q in 2, 3, 5, 7, 10 minutes

  \begin{itemize}
  \tightlist
  \item
    2 to 10 min lines up with the lag from the previous test
  \end{itemize}
\item
  One important clarification before the result

  \begin{itemize}
  \tightlist
  \item
    a martingale does not mean motionless
  \item
    spot moves plenty, with high volatility
  \item
    the test asks whether the next return is predictable from the past
    pattern of returns at this frequency
  \item
    a martingale says no
  \end{itemize}
\item
  The result on screen has two lines, one for each series
\item
  Spot:

  \begin{itemize}
  \tightlist
  \item
    VR(q) within 1\% of 1.00 at every horizon
  \item
    Hurst H = 0.498, CI contains 0.5
  \item
    spot satisfies the random-walk null at minute frequency
  \end{itemize}
\item
  P(UP):

  \begin{itemize}
  \tightlist
  \item
    VR(2) = 1.24, VR(10) = 1.41
  \item
    Hurst H = 0.537, CI excludes 0.5
  \item
    P(UP) does not satisfy the null
  \end{itemize}
\item
  Economically:

  \begin{itemize}
  \tightlist
  \item
    when P(UP) moves, more moves in the same direction tend to follow
    over the next several minutes
  \item
    P(UP) under-reacts to its own information
  \end{itemize}
\item
  Two timing results now, from independent angles, same answer

  \begin{itemize}
  \tightlist
  \item
    spot leads Kalshi by 2 min (the cross-series test)
  \item
    Kalshi under-reacts to itself at 2 to 10 min horizons (the
    within-series test)
  \end{itemize}
\item
  So the next question is practical

  \begin{itemize}
  \tightlist
  \item
    does that timing difference actually matter when you use Kalshi
    features to predict contract outcomes?
  \item
    or is it a microstructure curiosity with no consequence for
    prediction?
  \item
    which is the next slide
  \end{itemize}
\end{itemize}

\begin{center}\rule{0.5\linewidth}{0.5pt}\end{center}

\subsection{Slide 12: Kalshi Features Add No Content Above
Spot}\label{slide-12-kalshi-features-add-no-content-above-spot}

\textbf{\textasciitilde65 sec (method + result combined)}

\begin{itemize}
\tightlist
\item
  Third diagnostic: head-to-head classification on the same
  out-of-sample contracts
\item
  Three feature sets, three logistic regressions

  \begin{itemize}
  \tightlist
  \item
    spot-only: 8 features from BTC spot in the first 5 minutes (returns,
    realized volatility, volume)
  \item
    Kalshi-only: 11 features I built from the Kalshi panel

    \begin{itemize}
    \tightlist
    \item
      opening market state: opening P(UP), conviction spread, log volume
    \item
      early dynamics: 1-min and 3-min momentum, mean and std of P(UP)
      over minutes 0-4
    \item
      cross-asset block: matched concurrent ETH contract features
    \end{itemize}
  \item
    combined: union, 19 features
  \end{itemize}
\item
  Compare out-of-sample AUCs with the paired DeLong test

  \begin{itemize}
  \tightlist
  \item
    source on the slide footer with DOI link
  \end{itemize}
\item
  Quick justifications:

  \begin{itemize}
  \tightlist
  \item
    logistic regression because the question is about feature
    information, not model complexity
  \item
    5-fold time-series CV (train on past, test on future), no random
    shuffling in financial time series
  \item
    DeLong is the standard paired non-parametric test for comparing AUCs
    on shared data
  \end{itemize}
\item
  The result, on 1,595 out-of-sample BTC contracts:

  \begin{itemize}
  \tightlist
  \item
    Spot-only AUC = 0.796, CI {[}0.78, 0.82{]}, 73\% accuracy
  \item
    Kalshi-only AUC = 0.693

    \begin{itemize}
    \tightlist
    \item
      this matches the whitepaper's Extended LR exactly
    \item
      so the sanity check on the Kalshi feature set passes
    \end{itemize}
  \item
    Combined AUC = 0.760

    \begin{itemize}
    \tightlist
    \item
      worse than spot-only
    \item
      DeLong p \textless{} 0.001 for both gaps
    \end{itemize}
  \end{itemize}
\item
  The combined model being worse than spot-only is the key signature
  here

  \begin{itemize}
  \tightlist
  \item
    it is the standard pattern of an irrelevant feature block adding
    estimation variance to a finite sample
  \item
    the spot tape already contains the information; adding Kalshi only
    adds noise to fit
  \end{itemize}
\item
  The takeaway:

  \begin{itemize}
  \tightlist
  \item
    Kalshi features carry no incremental information above the spot tape
    in the same window
  \item
    coherent with the lead-lag finding (Kalshi is 2 min downstream)
  \item
    coherent with the variance-ratio finding (Kalshi under-reacts even
    to itself)
  \end{itemize}
\item
  Three diagnostics, same direction

  \begin{itemize}
  \tightlist
  \item
    Kalshi is delayed sentiment that is fully subsumed by spot for
    short-horizon directional prediction
  \end{itemize}
\end{itemize}

\begin{center}\rule{0.5\linewidth}{0.5pt}\end{center}

\subsection{Slide 13: Limitations and Next
Steps}\label{slide-13-limitations-and-next-steps}

\textbf{\textasciitilde50 sec}

\begin{itemize}
\tightlist
\item
  Affirmative finding first

  \begin{itemize}
  \tightlist
  \item
    the three diagnostics jointly characterize Kalshi as a specific kind
    of data source

    \begin{itemize}
    \tightlist
    \item
      calibrated
    \item
      real-money
    \item
      minute-frequency measurement of crowd belief
    \item
      that updates about 2 minutes after the spot tape
    \item
      and exhibits positive autocorrelation in its own changes
    \end{itemize}
  \item
    that is a description of what the data is, not a verdict that it has
    no value
  \end{itemize}
\item
  What an unconvinced referee would push on

  \begin{itemize}
  \tightlist
  \item
    Identification:

    \begin{itemize}
    \tightlist
    \item
      the 2-min lag confounds three sources

      \begin{itemize}
      \tightlist
      \item
        Coinbase candle aggregation
      \item
        scraper averaging of mid-quotes over the prior minute
      \item
        substantive reaction lag
      \end{itemize}
    \item
      I bound the substantive component at at least 1 full minute
    \item
      tick-level spot would identify the exact share
    \end{itemize}
  \item
    Inference:

    \begin{itemize}
    \tightlist
    \item
      cluster bootstrap by contract preserves within-contract dependence
    \item
      but treats contracts as independent draws across calendar time
    \item
      adjacent contracts share macro state
    \item
      a calendar-time block bootstrap would be more conservative
    \item
      the qualitative findings are robust but headline magnitudes could
      tighten
    \end{itemize}
  \item
    Generalizability:

    \begin{itemize}
    \tightlist
    \item
      6 weeks of moderate volatility is one regime
    \item
      the 2-min lag and the 0.10 AUC gap should not be treated as
      structural until they replicate during stress events (CPI
      releases, BTC-specific news shocks)
    \end{itemize}
  \end{itemize}
\item
  Two natural extensions I want to run

  \begin{itemize}
  \tightlist
  \item
    Hasbrouck 1995 information share decomposition

    \begin{itemize}
    \tightlist
    \item
      vector error correction model of Kalshi-implied price and spot
    \item
      the standard tool for venue-level price discovery
    \end{itemize}
  \item
    second-moment channel

    \begin{itemize}
    \tightlist
    \item
      whether the conviction spread (\textbar P(UP) - 0.5\textbar)
      forecasts realized volatility above HAR-RV and GARCH benchmarks
    \item
      the most promising remaining use of this data source
    \end{itemize}
  \end{itemize}
\end{itemize}

\begin{center}\rule{0.5\linewidth}{0.5pt}\end{center}

\subsection{Slide 15: Questions}\label{slide-15-questions}

\textbf{\textasciitilde15 sec}

\begin{itemize}
\tightlist
\item
  Three takeaways about Kalshi as a data source

  \begin{itemize}
  \tightlist
  \item
    Calibration:

    \begin{itemize}
    \tightlist
    \item
      unlike text-based sentiment proxies, P(UP) is
      information-theoretically meaningful
    \item
      probabilities, not orderings
    \end{itemize}
  \item
    Speed:

    \begin{itemize}
    \tightlist
    \item
      the 2-min lag against spot and the positive autocorrelation in
      P(UP)
    \item
      are measurements of how fast a real-money crowd absorbs price
      information at 1-minute granularity
    \end{itemize}
  \item
    Scope:

    \begin{itemize}
    \tightlist
    \item
      the directional channel is already in the spot tape
    \item
      so the use cases for this data source are second-moment
      forecasting and calibrated minute-frequency sentiment, not alpha
      in the underlying
    \end{itemize}
  \end{itemize}
\item
  Happy to take questions
\end{itemize}

\begin{center}\rule{0.5\linewidth}{0.5pt}\end{center}

\subsection{Backup Q\&A}\label{backup-qa}

\textbf{Q: Why join Kalshi to spot by exact UTC minute rather than
nearest timestamp?} - Exact UTC minute is conservative - the scraper
already averages mid-quotes over the prior minute - so the Kalshi value
at minute t is the average over (t-1, t{]} - and the Coinbase candle at
minute t is also (t-1, t{]} - Nearest-timestamp matching would smear the
join window and bias the contemporaneous correlation upward -
artifactually making PM look more contemporaneous with spot than it is

\textbf{Q: Why cluster bootstrap by contract rather than by day or
hour?} - Within-contract observations are mechanically dependent - P(UP)
at minute t in a contract is conditioned on P(UP) at minute t-1 in the
same contract - Contract is the smallest unit that is plausibly
exchangeable across calendar time - day or hour clustering would be more
conservative but coarser - I note this in the limitations slide

\textbf{Q: Why 5-min feature window specifically? Why not 3, or 7?} -
Bias-variance tradeoff - shorter window leaves more forward window to
predict on, but features are noisier - longer window gives cleaner
features but the contract has already partly converged - 5 minutes
balances both - the convergence curve Xavi showed is still in the early
plateau at t=5 - and the 10 remaining minutes is the prediction window
that matters for trading

\textbf{Q: Could the result reverse during high-volatility regimes?} -
Possible and worth testing - PM venues may be slower to update during
stress (wider spreads, thinner liquidity) - which would widen the lag,
not narrow it - But it is also possible that informed traders
concentrate on PM during news events - in which case PM could lead spot
transiently - Six weeks of moderate vol is one regime; this is the
generalizability point on slide 13

\textbf{Q: Why no XGBoost or nonlinear model in the head-to-head?} - The
question is about feature information, not model capacity - if Kalshi
features carried orthogonal information, a linear model would already
pick up part of it - and the combined-minus-spot delta would be
positive, even if small - The result is the opposite (combined is
worse), which is robust to model choice - a nonlinear model could in
principle widen the spot-Kalshi gap, but cannot reverse the sign of the
incremental delta in this finite sample

\end{document}
